\section*{Conclusion}
\addcontentsline{toc}{section}{Conclusion}

Ce devoir a permis d'analyser en détail le comportement d'un système de régulation de glycémie basé sur le modèle minimal de Bergman avec un correcteur PID.

La première partie a établi les bases analytiques du système. Le point d'équilibre en boucle ouverte, trouvé par \texttt{trim()}, donne $G_e = 46$ mg/dL, ce qui est en dehors de la zone de sécurité. En boucle fermée avec $K_p = 0.5$, $K_d = 0.01$ et $K_i = 0.1$, le point d'équilibre se déplace à $G_e = 92$ mg/dL, soit exactement le niveau basal $G_b$. La linéarisation autour de ce point a permis d'extraire les fonctions de transfert $G_{u_{BF}}$ et $G_{m_{BF}}$, dont les pôles complexes conjugués ($-0.0057 \pm 0.0545i$) expliquent le comportement oscillatoire observé dans les simulations. L'analyse de Routh-Hurwitz a montré que pour un correcteur purement proportionnel, la stabilité n'est garantie que pour $K_p < 0.008403$. Le lieu des racines avec le correcteur PID complet confirme la stabilité pour les gains choisis, et l'erreur en régime permanent est nulle grâce à l'action intégrale du contrôleur.

La deuxième partie a montré que la boucle ouverte est inadaptée: la glycémie quitte la zone de sécurité durant le transitoire et se stabilise à 82.5 mg/dL, en dessous de $G_b$. La boucle fermée ramène correctement la glycémie à $G_b = 92$ mg/dL malgré des oscillations transitoires, ce qui confirme l'intérêt du contrôleur PID pour ce type de système.

La troisième partie a mis en évidence les limites du modèle linéarisé. Pour des perturbations de faible amplitude, le modèle linéarisé capture qualitativement le comportement du système, mais prédit un régime transitoire plus oscillatoire que le modèle non linéaire. L'étude de l'effet de $K_p$ révèle que la re-linéarisation à chaque valeur de $K_p$ modifie les pôles du système, rendant la réponse moins oscillatoire à $K_p$ plus élevé dans le modèle linéarisé. Enfin, pour une perturbation dix fois plus grande ($m = 50$), le modèle linéarisé produit des prédictions physiologiquement impossibles (glycémie dépassant 950 mg/dL et passant sous zéro), confirmant que les hypothèses de linéarisation ne tiennent plus loin du point d'équilibre.

Ces résultats montrent qu'un contrôleur PID linéaire conçu autour d'un point d'équilibre fixe est insuffisant pour garantir la sécurité du patient face à des perturbations importantes. Une approche non linéaire ou une re-linéarisation adaptative serait nécessaire pour un système de pancréas artificiel robuste en conditions réelles.
